\documentclass[10pt,a4paper]{amsart}
\usepackage[margin=19mm]{geometry}
\usepackage{amsmath,amssymb,mathtools}
\usepackage[scr=rsfs,cal=euler]{mathalfa}
\usepackage{xfrac}
\usepackage[unicode,hidelinks,bookmarks=false]{hyperref}
\newcommand{\ie}{i.e.\ }
\newcommand{\cf}{cf.\ }
\newcommand{\resp}{respectively\ }
\newcommand{\Subsection}{\subsection}
\providecommand{\nobreakdash}{\nobreak\hbox{-}}
\setlength{\emergencystretch}{2em}
\multlinegap0pt
\usepackage{amssymb}
\usepackage{dsfont}
\usepackage{tikz}
\usepackage[normalem]{ulem}
\newtheorem{lemma}{Lemma}
\title{\vspace{-0.9cm} Two Erd\H{o}s--Hajnal-type theorems for forbidden order-size pairs}
\author{Fabian Arnold, Lior Gishboliner, Benny Sudakov}
\date{Source: arXiv:2406.04154v3 (CC BY 4.0)}
\begin{document}
\maketitle

\noindent\emph{Source and licence.} \vspace{-0.9cm} Two Erd\H{o}s--Hajnal-type theorems for forbidden order-size pairs. Version arXiv:2406.04154v3 (CC BY 4.0), distributed by arXiv under the Creative Commons Attribution 4.0 International licence, http://creativecommons.org/licenses/by/4.0/, as recorded in the arXiv licence metadata for this exact version and shown on the abstract page https://arxiv.org/abs/2406.04154v3. Full licence notice: \texttt{licenses/arxiv-2406.04154-CC-BY-4.0.txt}.\par\smallskip

\noindent\textit{Editorial scope.} Counted unit: the article's lemma general-values-lemma (source0.tex lines 1471--1481). The article's complete original proof of it is delivered with it (source0.tex lines 1489--1596, one \texttt{proof} environment of 108 lines, which the article places away from the statement). No mathematics is written, repaired or re-derived: the statement and the proof are the article's own lines, in the article's own notation, and no retained line is substituted or re-flowed.  No further source-local context is needed: every cross-reference of every retained line resolves inside the two delivered spans.  Nothing was omitted from any retained span, so the package carries no omission record.  No retained line contains a citation, so the wrapper carries no bibliography and no bibliography entry.  No retained line contains a figure environment, an image-inclusion command or drawing code, so no image is delivered and the package declares no asset dependency.  Editorial formatting only, in addition: an AMS \texttt{amsart} preamble that restates the article's own declarations and macros used by the retained lines, namely the environments \texttt{lemma} on separate counters, together with the standard packages those lines need.  The retained environments print in this document as Lemma 1 in order of appearance, whereas the article prints the same items with the numbering of its own full document.  The complete native source of the article as distributed in the arXiv e-print for this exact version is bundled unchanged as \texttt{sources/160406/source0.tex}.
\begin{lemma} \label{general-values-lemma}
    Let $A,B,C,D,E \in \mathbb{Q}$ such that $B \neq 0$ or $A,C \neq 0$.
    Then the function 
    % $f: X = \{(x_1, \dots, x_m) \in \mathbb{Z}_{\geq 0}^t | x_1 + \dots + x_m = m\} \to \mathbb{Q}$ with
	$$
	f(x_1,\dots,x_m) = (Am+D) \sum_{1 \leq i \leq m} x_i^2 + 
	B \sum_{1 \leq i \leq m} x_i^3 + C \sum_{1 \leq i < j \leq m} x_i x_j (x_i - x_j)+E
	$$
    takes $\Omega(m^3)$ distinct values on the set of integer vectors $(x_1,\dots,x_m)$ with $x_1,\dots,x_m \geq 0$ and 
    $\sum_{i=1}^m x_i = m$.
\end{lemma}

\begin{proof}[Proof of Lemma \ref{general-values-lemma}]
    We may and will assume that $m$ is large enough as a function of $A,B,C,D$. 
    We first consider ``symmetric" vectors $x = (x_1,\dots,x_t)$, i.e., vectors satisfying $x_i = x_{t+1-i}$ for every $1 \leq i \leq t$. 
    To keep the notation cleaner, we ignore (i.e., do not write) coordinates $i \in [m]$ with $x_i = 0$, as these coordinates do not contribute to $f$. Namely, when we write $(x_1,\dots,x_t)$, we actually mean $(x_1,\dots,x_m)$ with $x_i = 0$ for $t < i \leq m$. 
    If $x = (x_1,\dots,x_t)$ is symmetric then 
    $$\sum_{1 \leq i < j \leq t} x_i x_j (x_i - x_j) = 
    \sum_{1 \leq i < j \leq t} x_i^2x_j - \sum_{1 \leq i < j \leq t} x_ix_j^2 = 0,$$ because $x_i^2x_j = x_{t+1-j}x_{t+1-i}^2$. This means that for symmetric vectors $x$, it holds that
    $$
    f(x) = (Am+ \nolinebreak D) \sum_{i} x_i^2 + 
	B \sum_{i} x_i^3 + E.
    $$
    
    
    Fix a constant $0 < \epsilon < \frac{1}{16}$ to be chosen later.
    For an integer $\epsilon m \leq \ell \leq 2\epsilon m$, let $x^{(\ell)}$ be the symmetric vector
    $$x^{(\ell)} = (\ell, \lfloor \sqrt{m} \rfloor, \dots, \lfloor \sqrt{m} \rfloor, 2, \dots, 2, 1, \dots, 1, 2, \dots, 2, \lfloor \sqrt{m} \rfloor, \dots, \lfloor \sqrt{m} \rfloor, \ell),$$ 
    where the number of coordinates equal to $\lfloor \sqrt{m} \rfloor$ is exactly $2 \lfloor \frac{\sqrt{m}}{8} \rfloor$; the number $2$s is exactly $2 \lfloor \frac{m}{16} \rfloor$; and the number of $1$s is chosen so that the sum of all coordinates is exactly $m$. 
    We have 
    $$
    2\ell+ \lfloor \sqrt{m} \rfloor \cdot 
    2 \left\lfloor \frac{\sqrt{m}}{8} \right\rfloor + 
    2 \cdot 2 \left\lfloor \frac{m}{16} \right\rfloor \leq 2\ell + \frac{m}{2} \leq 4\epsilon m + \frac{m}{2} \leq \frac{3m}{4},
    $$
    so the number of $1$s is at least $\frac{m}{4}$. 
    
    % At first, we will consider "symmetric" vectors in $X$, i.e. $x_1 = x_m, x_2 = x_{t-1}, \dots$ (this ensures that the last term of $f$ is zero), and will neglect writing out the coordinates where $x_i = 0$ (which do not contribute to $f$). 
    % We start by considering the vectors of the form $x^{(l)} = (l, \lfloor \sqrt{m} \rfloor, \dots, \lfloor \sqrt{m} \rfloor, 2, \dots, 2, 1, \dots, $ $1, 2, \dots, 2, \lfloor \sqrt{m} \rfloor, \dots, \lfloor \sqrt{m} \rfloor, l) \in X$ for $\epsilon m \leq l \leq 2\epsilon m$, where $0 < \epsilon < \frac{1}{8}$ (which we will specify below) and we fix exactly $2 \lfloor \frac{\sqrt{m}}{8} \rfloor$ $\lfloor \sqrt{m} \rfloor$'s and $2 \lfloor \frac{m}{16} \rfloor$ $2$'s (the amount of $1$'s varies but is at least $\frac{m}{4}$). 

    By the assumptions of the lemma, $A,B$ cannot both be $0$. 
    For convenience, we will assume from now on that $A \geq 0$, and if $A = 0$ then $B > 0$ (this can be guaranteed by multiplying $f$ by $-1$, if necessary).
    We claim that $f(x^{(\ell+1)}) - f(x^{(\ell)}) = \Omega(m^2)$. 
    % Note that $x^{(l+1)}$ is obtained from $x^{(l)}$ by increasing $l$ to $l+1$ and decreasing the number of $1$s by two. 
    Indeed, recall that $x^{(\ell)}$ and $x^{(\ell+1)}$ have the same number of coordinates which equal $2$ and the same number of coordinates which equal $\lfloor \sqrt{m} \rfloor$. Also, the term $\sum_{1 \leq i < j \leq m} x_i x_j (x_i - x_j)$ in the definition of $f$ vanishes on both $x^{(\ell)}$ and $x^{(\ell+1)}$. Hence,
    \begin{align*}
        f(x^{(\ell+1)}) - f(x^{(\ell)})   
        &= 2(Am+D)\cdot ((\ell+1)^2 - \ell^2 - 1^2) + 2B \cdot ((\ell+1)^3 - \ell^3 - 1^3) \\
        &= 2\ell \cdot (2Am + 3B\ell + 2D + 3B).
    \end{align*}
    If $A,B \neq 0$ then choose $\epsilon$ to satisfy $\epsilon \leq \frac{A}{12|B|}$ (if $A$ or $B$ is 0 then there is no restriction).
    This ensures that if $A \neq 0$ then
    $|3B\ell| \leq 3|B| \cdot 2 \epsilon m \leq \frac{1}{2}Am$. Together with 
    $|2D + 3B| \leq \frac{1}{2}Am$ for $m$ large enough, we get that
    $2Am + 3B\ell + 2D + 3B \geq Am$ and thus 
    $f(x^{(\ell+1)}) - f(x^{(\ell)}) = 2\ell \cdot (2Am + 3B\ell + 2D + 3B) \geq 2 \epsilon m \cdot Am = \Omega(m^2)$, as required.
    If $A = 0$, then by assumption $B > 0$. Also, $|2D + 3B| \leq B \epsilon m \leq B\ell$ for $m$ large enough. Hence, $2\ell \cdot(2Am + 3B\ell + 2D + 3B) \geq 2\ell \cdot 2B\ell \geq
    2 \epsilon m \cdot 2B \epsilon m = 4B \epsilon^2 m^2 = \Omega(m^2)$.
    This proves our claim that $f(x^{(\ell+1)}) - f(x^{(\ell)}) = \Omega(m^2)$.
    % Note that in both cases, the sign of the difference is equal to the sign of $A$ (or $B$) and does not depend on $l$. Thus we get $\Omega(m)$ different values of $f$, each bigger (or each smaller) than the previous one by $\Omega(m^2)$.

    So far, we showed that there exist $\Omega(m)$ vectors $x^{(\ell)}$, $\epsilon m \leq \ell \leq 2\epsilon m$, such that the difference in the $f$-values of two consecutive vectors is 
    $\Omega(m^2)$. 
	For the rest of the proof, we fix any $\epsilon m \leq \ell \leq 2\epsilon m$ 
    and perform certain operations on $x^{(\ell)}$ to get vectors $x = (x_1,\dots,x_m)$ with 
    $\Omega(m^2)$ different $f$-values which all lie between $f(x^{(\ell)})$ and $f(x^{(\ell+1)})$ or between $f(x^{(\ell-1)})$ and $f(x^{(\ell)})$. Doing this for every $\ell$ gives us $\Omega(m^3)$ different $f$-values, establishing the lemma. We consider several cases.

    \paragraph{Case 1:} $C \neq 0$.  
    In this case, starting with the vector $x^{(\ell)}$, we repeatedly apply the following operation: take a coordinate $i$ with $x_i = 1, x_{i+1} = 2$, and change this to $x_i = 2, x_{i+1} = 1$ (note that the resulting vectors are no longer symmetric). 
    One such operation changes the value of $f$ by 
    $C(2 \cdot 1(2-1) - 1 \cdot 2(1-2)) = 4C$. 
    Also, it is possible to perform this operation as long as there is a coordinate $i$ with $x_i = 1, x_{i+1} = 2$. Initially (i.e., in $x^{(\ell)}$), there are at least $\frac{m}{4}$ $1$s which come directly before $\lfloor \frac{m}{16} \rfloor$ $2$s, and we can perform the operation until we reach the situation where these $2$s come before the $1$s. Hence, we can perform this operation at least 
    $\frac{m}{4} \lfloor \frac{m}{16} \rfloor = \Omega(m^2)$ times. Therefore, there are $\Omega(m^2)$ different values of $f$, with two consecutive values at distance exactly $4C = \Theta(1)$. This means that there are $\Omega(m^2)$ different $f$-values between $f(x^{(\ell)})$ and $f(x^{(\ell+1)})$ or between 
    $f(x^{(\ell-1)})$ and $f(x^{(\ell)})$ (depending on the sign of $C$), as required. 
    
    \paragraph{Case 2:} $C = 0$. 
    So $f(x) = (Am+ \nolinebreak D) \sum_{i} x_i^2 + 
	B \sum_{i} x_i^3 + E$.
    Also, by assumption, $B \neq 0$. 
    In this case we will use two operations:
    \begin{itemize}
        \item An operation which changes the value of $f$ by some fixed number $p = \Theta(1)$.
        \item An operation which changes the value of $f$ by some number $q$ with $\alpha m \leq q \leq \beta m$, for some fixed absolute constants $\alpha,\beta > 0$.
    \end{itemize}
    Moreover, we will be able to perform each of these operations $\Omega(m)$ times independently of each other (again starting from $x^{(\ell)}$). 
    % This will imply that there are $\Omega(m^2)$ different values of $f$ between $f(x^{(\ell)})$ and $f(x^{(\ell+1)})$ or between $f(x^{(\ell-1)})$ and $f(x^{(\ell)})$ (depending on the sign of $q$), 
    This will imply that there are $\Omega(m^2)$ different values of $f$ between $f(x^{(\ell)})$ and $f(x^{(\ell+1)})$, 
    because by applying the second operation $\Omega(m)$ times we get $\Omega(m)$ $f$-values which are $\Theta(m)$ apart, and we can then ``fill in these gaps" by applying the first operation $\Omega(m)$ times.
    
    The first operation is as follows: Choose distinct coordinates $i,j,k,h$ with $x_i = x_j = x_k = 2, x_h = 0$, and change this to $x_i = x_j = x_k = 1, x_h = 3$. This changes $f$ by 
    $$(Am+D)(3 \cdot 1^2 + 3^2 - 3 \cdot 2^2) + B(3 \cdot 1^3 + 3^3 - 3 \cdot 2^3) = 6B =: p = \Theta(1).$$ 
    Recall that $x^{(\ell)}$ has exactly $2 \cdot \lfloor \frac{m}{16} \rfloor$ coordinates which equal $2$. Since the sum of all coordinates is exactly $m$, and the total number of coordinates is $m$, there are also at least $2 \cdot \lfloor \frac{m}{16} \rfloor$ coordinates which equal $0$. 
    Therefore, it is possible to perform this operation at least 
    $\lfloor \frac{2}{3} \cdot \lfloor \frac{m}{16} \rfloor \rfloor = \Omega(m)$ times. 
    
    It remains to describe the second operation. Here we distinguish between two cases: 
    % For each value of $A$, we will also find an operation changing $f$ by $\Theta(m)$ that we can perform independently from the first operation:
    \paragraph{Case 2.1:} $A \neq 0$, so $A > 0$ (by our assumption). In this case we use the following operation: 
    Take coordinates $i,j$ with $x_i = x_j = 1$, and change this to $x_i = 2, x_j = 0$. This changes $f$ by 
    $
    q := (Am+D)(2^2 - 2 \cdot 1^2) + B(2^3 - 2 \cdot 1^3) = 2(Am+D) + 6B.
    $
    For large enough $m$ we have $Am \leq q \leq 3Am$, so $q = \Theta(m)$. 
    Since $x^{(\ell)}$ has at least $\frac{m}{4}$ $1$s, it is possible to perform this operation at least $\lfloor \frac{m}{8} \rfloor = \Omega(m)$ times, regardless of the number of times that the first operation was performed. 
    
    \paragraph{Case 2.2:} $A = 0$, so $B > 0$.
    In this case we use the following operation: take coordinates $i,j$ with  
    $\lfloor \sqrt{m} \rfloor \leq x_i \leq 2 \lfloor \sqrt{m} \rfloor$ and $x_j = 1$, increase $x_i$ by $1$, and set $x_j = 0$. This changes $f$ by 
    % $$q := D((x_i+1)^2 - x_i^2 - 1^2) + B((x_i+1)^3 - x_i^3 - 1^3) = x_i (2D + 3B(x_i+1))  \geq 2Bx_i^2 = \Theta(x_i^2) = \Theta(m),$$
    $$q := D((x_i+1)^2 - x_i^2 - 1^2) + B((x_i+1)^3 - x_i^3 - 1^3) = 
    3Bx_i^2 + (2D+3B)x_i.
    $$
    As $x_i = \Theta(\sqrt{m})$, we have $q = \Theta(m)$. 
    Recall that $x^{(\ell)}$ has exactly $2 \lfloor \frac{\sqrt{m}}{8} \rfloor$ coordinates which equal $\lfloor \sqrt{m} \rfloor$. For each such coordinate $i$, it is possible to perform this operation $\lfloor \sqrt{m} \rfloor$ times on $x_i$ (until $x_i$ becomes larger than $2 \lfloor \sqrt{m} \rfloor$). Also, $x^{(\ell)}$ has at least $\frac{m}{4}$ $1$s, meaning that there are enough $1$s for all of these $\Omega(m) \leq 2 \lfloor \frac{\sqrt{m}}{8} \rfloor \cdot \lfloor \sqrt{m} \rfloor \leq \frac{m}{4}$ operations. 
    So we can perform this second operation $\Omega(m)$ times, regardless of how many times we used the first operation.
    This completes the proof of the lemma. 
    % It is possible to perform this operation at least $\lfloor \sqrt{m} \rfloor 2 \lfloor \frac{\sqrt{m}}{8} \rfloor = \Omega(m)$ times. 
    %
    % In each case, we have found $\Theta(m^2)$ different values of $f$ between $f(x^{(l)})$ and $f(x^{(l+1)})$, so we found a total of $\Omega(m^3)$ distinct values of $f$.
\end{proof}

\end{document}
